\chapter{भिन्नें और घातें}\label{ch:g9:fractions}

भिन्नें और घातें रोज़ की गणना के औज़ार हैं: राशियाँ बाँटना, अनुपात
तुलना करना, बहुत बड़ी या बहुत छोटी संख्याएँ लिखना। यह अध्याय उनके
साथ गणना करने के नियम पक्के करता है --- और हर बीच का क़दम लिखकर ---
तथा \omterm{def:g9:fractions:scientific}{वैज्ञानिक लेखन} से परिचय कराता है।

\section{भिन्नों के साथ गणना}

\begin{definition}[भिन्न]\label{def:g9:fractions:fraction}
\emph{भिन्न}\index{भिन्न} $\dfrac ab$ ($b \neq 0$ के साथ) $a$ में
$b$ का \omterm{def:g3:division:remainder}{भागफल} है। दो भिन्नें तब बराबर होती हैं जब एक को दूसरी से,
\omterm{def:g4:fractions:def}{अंश} और हर दोनों को एक ही शून्येतर संख्या से गुणा (या भाग) करके पाया
जा सके:
\[
\frac ab = \frac{a \times k}{b \times k} \qquad (k \neq 0).
\]
\end{definition}

\begin{example}\label{ex:g9:fractions:simplify}
$\dfrac{42}{56}$ को क़दम दर क़दम सरल करो:
\[
\frac{42}{56} = \frac{42 \div 2}{56 \div 2} = \frac{21}{28}
= \frac{21 \div 7}{28 \div 7} = \frac34 .
\]
(\cref{ch:g9:arith} में महत्तम समापवर्तक यही काम एक ही क़दम में कर
देगा।)
\end{example}

\begin{method}[भिन्नें जोड़ना और घटाना]\label{met:g9:fractions:add}
\begin{enumerate}
  \item दोनों भिन्नों को एक \emph{उभयनिष्ठ हर} पर लाओ (यानी दोनों
        हरों का कोई उभयनिष्ठ \omterm{def:g5:division:multiple}{गुणज});
  \item \omterm{def:g4:fractions:def}{अंश} जोड़ो या घटाओ, और हर वही रहने दो;
  \item हो सके तो उत्तर को सरल करो।
\end{enumerate}
\end{method}

\begin{example}\label{ex:g9:fractions:add}
$\dfrac56 + \dfrac34$ निकालो। $6$ और $4$ का एक उभयनिष्ठ हर $12$
है:
\[
\frac56 + \frac34
= \frac{5 \times 2}{6 \times 2} + \frac{3 \times 3}{4 \times 3}
= \frac{10}{12} + \frac{9}{12}
= \frac{19}{12}.
\]
$2 - \dfrac37$ निकालो। $2$ को $7$ के ऊपर \omterm{def:g9:fractions:fraction}{भिन्न} की तरह लिखो:
\[
2 - \frac37 = \frac{14}{7} - \frac37 = \frac{11}{7}.
\]
\end{example}

\begin{omfigure}
\begin{tikzpicture}[scale=0.55]
  \foreach \i in {0,...,11}
    \draw[omExo] ({mod(\i,12)},0) rectangle ({mod(\i,12)+1},1);
  \foreach \i in {0,...,9}
    \fill[omDef!35] (\i,0) rectangle (\i+1,1);
  \draw[thick] (0,0) rectangle (12,1);
  \node[below, font=\small] at (6,-0.1)
    {$\frac56 = \frac{10}{12}$: 12 में से 10 ख़ाने};
  \begin{scope}[yshift=-2.4cm]
  \foreach \i in {0,...,11}
    \draw[omExo] (\i,0) rectangle (\i+1,1);
  \foreach \i in {0,...,8}
    \fill[omThm!35] (\i,0) rectangle (\i+1,1);
  \draw[thick] (0,0) rectangle (12,1);
  \node[below, font=\small] at (6,-0.1)
    {$\frac34 = \frac{9}{12}$: 12 में से 9 ख़ाने};
  \end{scope}
\end{tikzpicture}

{\small उभयनिष्ठ हर क्यों: दोनों पूरों को $12$ बराबर ख़ानों में काट
देने पर भिन्नें तुलने और जुड़ने लायक़ हो जाती हैं --- दोनों मिलकर
$10 + 9 = 19$ बारहवें भाग।}
\end{omfigure}

\begin{method}[भिन्नों का गुणा और भाग]\label{met:g9:fractions:mult}
\begin{enumerate}
  \item गुणा करने के लिए \omterm{def:g4:fractions:def}{अंश} आपस में और हर आपस में गुणा करो:
        $\dfrac ab \times \dfrac cd = \dfrac{ac}{bd}$ (जहाँ हो सके,
        गुणा करने से \emph{पहले} सरल कर लो);
  \item भाग देने के लिए \omterm{def:g5:division:multiple}{भाजक} के \omterm{def:g8:fractions:inverse}{व्युत्क्रम} से गुणा करो:
        $\dfrac ab \div \dfrac cd = \dfrac ab \times \dfrac dc$।
\end{enumerate}
\end{method}

\begin{example}\label{ex:g9:fractions:mult}
\[
\frac{7}{15} \times \frac{25}{14}
= \frac{7 \times 25}{15 \times 14}
= \frac{7 \times 25}{15 \times 2 \times 7}
= \frac{25}{30} = \frac56,
\]
जहाँ हमने पहले उभयनिष्ठ गुणनखंड $7$ से, फिर $5$ से सरल किया। और अब
एक भाग:
\[
\frac{3}{4} \div \frac{9}{8}
= \frac34 \times \frac89
= \frac{3 \times 8}{4 \times 9}
= \frac{24}{36} = \frac23 .
\]
\end{example}

\section{घातें}

\begin{definition}[घात]\label{def:g9:fractions:power}
किसी वास्तविक संख्या $a$ और किसी धन पूर्णांक $n$ के लिए
\emph{घात}\index{घात} $a^n$, $a$ के बराबर $n$ गुणकों का \omterm{def:g2:mult:def}{गुणनफल} है:
\[
a^n = \underbrace{a \times a \times \dots \times a}_{n \text{ गुणक}} .
\]
यह तय किया गया है कि $a^0 = 1$ ($a \neq 0$ के लिए), और ऋण \omterm{def:g8:powers:def}{घातांक}
\omterm{def:g8:fractions:inverse}{व्युत्क्रम} दिखाते हैं:
\[
a^{-n} = \frac{1}{a^n}.
\]
\end{definition}

\begin{example}
$2^4 = 2 \times 2 \times 2 \times 2 = 16$; \quad
$(-3)^2 = 9$ पर $-3^2 = -(3^2) = -9$ (\omterm{def:g8:powers:def}{घातांक} ऋण चिह्न से पहले लगता
है); \quad $10^{-3} = \frac{1}{1000} = 0.001$; \quad
$5^1 = 5$ और $5^0 = 1$।
\end{example}

\begin{theorem}[घातांक के नियम]\label{thm:g9:fractions:rules}
सभी शून्येतर $a$, $b$ और सभी पूर्णांकों $m$, $n$ के लिए:
\[
a^m \times a^n = a^{m+n},
\qquad
\frac{a^m}{a^n} = a^{m-n},
\qquad
\left(a^m\right)^n = a^{mn},
\qquad
(ab)^n = a^n b^n .
\]
\end{theorem}

\begin{proof}
धन घातांकों के लिए गुणक गिन लो। पहले नियम के लिए:
$a^m \times a^n$ में $a$ के $m$ गुणकों के बाद $a$ के $n$ गुणक आते
हैं, यानी कुल $m + n$ गुणक। तीसरे के लिए:
$\left(a^m\right)^n$ में $m$ गुणकों का गुच्छा $n$ बार दोहराया जाता
है, यानी $mn$ गुणक। बाक़ी दोनों नियम भी वैसे ही हैं; और \omterm{def:g1:counting:zero}{शून्य} तथा ऋण
घातांकों तक बढ़ाव $a^{-n} = \frac{1}{a^n}$ से जाँच लिया जाता है (और
इसी से $a^0 = 1$ अकेला निभने वाला चुनाव बन जाता है, क्योंकि
$a^n \times a^0$ को $a^{n+0} = a^n$ होना ही है)।
\end{proof}

\begin{example}\label{ex:g9:fractions:rules}
क़दम दर क़दम सरल करो:
\[
\frac{3^7 \times 3^{-2}}{3^4}
= \frac{3^{7 + (-2)}}{3^4}
= \frac{3^5}{3^4}
= 3^{5-4} = 3^1 = 3 .
\]
और दो अक्षरों के साथ: $(2a^3)^2 \times a = 4 a^6 \times a = 4a^7$।
\end{example}

\section{दस की घातें और वैज्ञानिक लेखन}

\begin{proposition}[दस की घातें]\label{prop:g9:fractions:ten}
किसी धन पूर्णांक $n$ के लिए: $10^n$ को ``$1$ के बाद $n$ \omterm{def:g1:counting:zero}{शून्य}''
लिखा जाता है, और $10^{-n} = 0.0\dots01$ होता है जिसमें $1$ दशमलव के
$n$-वें स्थान पर बैठता है। किसी \omterm{def:g6:decimals:places}{दशमलव संख्या} को $10^n$ (या
$10^{-n}$) से गुणा करने पर उसका \omterm{def:g4:decimals:point}{दशमलव बिंदु} $n$ स्थान दाईं (या
बाईं) ओर सरक जाता है।
\end{proposition}

\begin{definition}[वैज्ञानिक लेखन]\label{def:g9:fractions:scientific}
किसी धन \omterm{def:g6:decimals:places}{दशमलव संख्या} का \emph{वैज्ञानिक
लेखन}\index{वैज्ञानिक लेखन} उसे इस रूप में लिखने का अकेला तरीक़ा है
\[
a \times 10^n,
\qquad\text{जहाँ } 1 \leq a < 10 \text{ और } n \text{ एक पूर्णांक।}
\]
\end{definition}

\begin{example}\label{ex:g9:fractions:scientific}
पृथ्वी--सूर्य की दूरी क़रीब $149\,600\,000$ km $= 1.496 \times
10^8$ km है। पानी के एक अणु का आकार क़रीब $0.000\,000\,000\,28$ m
$= 2.8 \times 10^{-10}$ m है। ऐसी संख्याओं के साथ गणना करने के लिए
दशमलव वाले हिस्से आपस में और दस की घातें आपस में जमा लो:
\[
(3 \times 10^5) \times (2.5 \times 10^{-2})
= (3 \times 2.5) \times 10^{5 + (-2)}
= 7.5 \times 10^{3}.
\]
\end{example}

\begin{method}[किसी संख्या को वैज्ञानिक लेखन में लिखना]\label{met:g9:fractions:scientific}
\begin{enumerate}
  \item \omterm{def:g4:decimals:point}{दशमलव बिंदु} को पहले शून्येतर अंक के ठीक बाद सरका दो; इससे
        गुणक $a$ मिल जाता है, जिसके लिए $1 \leq a < 10$;
  \item गिनो कि बिंदु कितने स्थान सरका: वही \omterm{def:g7:stats:frequency}{गिनती} $n$ है, और वह धन
        होती है अगर असली संख्या $\geq 10$ थी, तथा ऋण अगर वह $< 1$
        थी;
  \item जाँचो: $a \times 10^n$ से असली संख्या वापस बननी चाहिए।
\end{enumerate}
\end{method}

\section{अभ्यास}

\begin{exercise}[$\star$]\label{exo:g9:fractions:1}
निकालो और उत्तर पूरी तरह सरल की हुई \omterm{def:g9:fractions:fraction}{भिन्न} के रूप में दो:
\[
\frac13 + \frac15, \qquad
\frac72 - \frac53, \qquad
\frac{5}{12} + \frac38, \qquad
3 - \frac45 .
\]
\end{exercise}

\begin{exercise}[$\star$]\label{exo:g9:fractions:2}
निकालो और सरल करो:
\[
\frac59 \times \frac{27}{35}, \qquad
\frac{8}{15} \div \frac{4}{5}, \qquad
\frac23 \times \frac34 \times \frac45 .
\]
\end{exercise}

\begin{exercise}[$\star$]\label{exo:g9:fractions:3}
गणना करो:
\[
2^5, \qquad (-2)^4, \qquad -2^4, \qquad 4^{-2}, \qquad
\left(\tfrac23\right)^3, \qquad 7^0 .
\]
\end{exercise}

\begin{exercise}[$\star$]\label{exo:g9:fractions:4}
एक ही \omterm{def:g9:fractions:power}{घात} के रूप में लिखो:
\[
5^3 \times 5^6, \qquad
\frac{7^9}{7^5}, \qquad
\left(2^3\right)^4, \qquad
\frac{3^2 \times 3^7}{3^5}, \qquad
2^5 \times 4 .
\]
\end{exercise}

\begin{exercise}[$\star$]\label{exo:g9:fractions:5}
\omterm{def:g9:fractions:scientific}{वैज्ञानिक लेखन} में लिखो: $45\,000$; $0.0072$; $302.5$;
$0.000\,000\,9$; एक करोड़ बीस लाख।
\end{exercise}

\begin{exercise}[$\star\star$]\label{exo:g9:fractions:6}
निकालो, और उत्तर \omterm{def:g9:fractions:scientific}{वैज्ञानिक लेखन} में दो:
\[
(2 \times 10^7) \times (4.5 \times 10^{-3}),
\qquad
\frac{9 \times 10^5}{3 \times 10^{-2}},
\qquad
(5 \times 10^3)^2 .
\]
\end{exercise}

\begin{exercise}[$\star\star$]\label{exo:g9:fractions:7}
एक टंकी $\frac35$ भरी है। $30$ लीटर और डालने पर वह $\frac34$ भर
जाती है। टंकी की \omterm{def:g2:measure:units}{धारिता} क्या है? (पहले टंकी का वह भाग निकालो जो
डाला गया।)
\end{exercise}

\begin{exercise}[$\star\star$]\label{exo:g9:fractions:8}
ऐलिस एक केक का $\frac14$ खाती है, फिर बेन जो बचा उसका $\frac13$
खाता है, फिर कार्ला जो अब भी बचा उसका $\frac12$। आख़िर में केक का
कितना भाग बचा? (हर क़दम के बाद बचा हुआ हिस्सा निकालो।)
\end{exercise}

\begin{exercise}[$\star\star$]\label{exo:g9:fractions:9}
प्रकाश हर सेकंड $3 \times 10^8$ मीटर चलता है। सूर्य क़रीब
$1.5 \times 10^{11}$ m दूर है। सूरज की रोशनी हम तक पहुँचने में
कितना समय लेती है? उत्तर सेकंड में दो (\omterm{def:g9:fractions:scientific}{वैज्ञानिक लेखन} की ज़रूरत
नहीं), फिर मिनट और सेकंड में।
\end{exercise}

\begin{exercise}[$\star\star\star$]\label{exo:g9:fractions:10}
कौन बड़ा है, $2^{30}$ या $3^{20}$? (संकेत:
$\left(a^m\right)^n = a^{mn}$ से दोनों को \omterm{def:g8:powers:def}{घातांक} $10$ वाली घातों के
रूप में लिखो, और $2^3$ की $3^2$ से तुलना करो।)
\end{exercise}

\section{समस्या: आवर्ती दशमलवों का गुप्त कूट}

\begin{problem}[{सप्ताहांत समस्या --- हर भिन्न का दशमलव लेखन या तो
रुक जाता है या दोहराता है, हर आवर्ती दशमलव एक भिन्न है, और
$0.999\ldots$ ठीक $1$ के बराबर है}]
\label{pb:g9:fractions:1}
$3$ में $11$ का भाग दो और अंक एक जाप में बदल जाते हैं:
$0.272727\ldots$ हमेशा-हमेशा। और $1$ में $4$ का भाग दो तो अंक ठप
हो जाते हैं: $0.25$। यह समस्या सिद्ध करती है कि किसी \omterm{def:g9:fractions:fraction}{भिन्न} के लिए
यही दो बरताव संभव हैं --- और फिर मशीन उलट देती है: कोई भी आवर्ती
दशमलव दो, वह उसके पीछे छिपी \omterm{def:g9:fractions:fraction}{भिन्न} वापस बना देती है। रास्ते में वह
स्कूली गणित की सबसे झगड़ालू बराबरी का फ़ैसला भी हमेशा के लिए कर देती
है।

\medskip
\noindent\textbf{भाग I --- \omterm{def:g9:fractions:fraction}{भिन्न} से दशमलव तक।}
\begin{enumerate}
  \item लंबे भाग से $\frac14$, $\frac13$, $\frac56$ और
        $\frac{3}{11}$ के दशमलव लेखन निकालो। कौन रुक जाते हैं,
        कौन दोहराते हैं, और दोहराने वाला गुच्छा क्या है?
  \item समझाओ कि \emph{किसी भी} \omterm{def:g9:fractions:fraction}{भिन्न} $\frac ab$ का दशमलव लेखन या
        तो रुकना क्यों पड़ेगा या आख़िरकार दोहराना। ($b$ से लंबे भाग
        में लगातार \omterm{def:g3:division:remainder}{शेषफल} कौन-कौन से मान ले सकते हैं? और जिस पल कोई
        \omterm{def:g3:division:remainder}{शेषफल} दूसरी बार आता है, तब क्या होता है?)
  \item $1$ में $7$ का भाग इतनी दूर तक ले जाओ कि दोहराने वाला
        गुच्छा मिल जाए। वह कितना लंबा है, और उस लंबाई की तुलना
        सवाल 2 में वादा की गई सीमा से कैसी है?
  \item कुछ भिन्नें इसलिए रुक जाती हैं कि उनका हर बढ़ाकर दस की \omterm{def:g9:fractions:power}{घात}
        बनाया जा सकता है: $\frac14 = \frac{25}{100}$,
        $\frac{3}{8} = \frac{375}{1000}$। यह कसौटी समझाओ, और यह भी
        कि कोई भी पूर्ण-संख्या गुणक $3$, $6$ या $11$ को $10$,
        $100$, $1000, \dots$ में कभी क्यों नहीं बदल सकता (\omterm{def:g2:mult:def}{गुणनफल}
        किस अंक पर ख़त्म होता?)।
  \item बिना भाग दिए $\frac{9}{40}$, $\frac{5}{12}$ और
        $\frac{33}{50}$ को ``रुकती हैं'' तथा ``दोहराती हैं'' में
        बाँटो, फिर जो दोनों रुकती हैं उनका दशमलव लेखन बताओ।
\end{enumerate}

\medskip
\noindent\textbf{भाग II --- आवर्ती दशमलव से वापस \omterm{def:g9:fractions:fraction}{भिन्न} तक।}
\begin{enumerate}[resume]
  \item सरकाने वाली तरकीब। मान लो $x = 0.7777\ldots$ है। $10x$
        निकालो, फिर $10x - x$, और इससे $x$ \omterm{def:g9:fractions:fraction}{भिन्न} के रूप में
        निकालो। भाग देकर जाँचो।
  \item मान लो $x = 0.727272\ldots$ है। दस की कौन-सी \omterm{def:g9:fractions:power}{घात} $x$ को
        ठीक एक दोहराने वाले गुच्छे जितना सरकाती है? उसी से $x$ को
        सबसे सरल \omterm{def:g9:fractions:fraction}{भिन्न} के रूप में लिखो।
  \item एक मिली-जुली स्थिति: $x = 0.58333\ldots$ (सिर्फ़ $3$
        दोहराता है)। $1000x - 100x$ निकालो और इससे $x$ को सबसे
        सरल \omterm{def:g9:fractions:fraction}{भिन्न} के रूप में निकालो।
  \item और वह मशहूर वाला: सरकाने वाली तरकीब
        $x = 0.9999\ldots$ पर लगाओ। तुम्हें कौन-सी \omterm{def:g9:fractions:fraction}{भिन्न} ---
        कौन-सी \emph{संख्या} --- मिलती है? इस नतीजे को अपनी समझ से
        मिलाओ, और \cref{pb:g6:decimals:1} वाले भूतिया पड़ोसी तथा
        \cref{pb:g6:fractions:1} वाले सिकुड़ते शेष को याद रखो: क्या
        $0.999\ldots$, $1$ से ``ज़रा-सा नीचे'' है, या $1$ का ही
        दूसरा नाम?
  \item $2.454545\ldots$ और $0.123123123\ldots$ को सबसे सरल
        भिन्नों में बदलो।
\end{enumerate}

\medskip
\noindent\textbf{भाग III --- शब्दकोश पूरा।}
\begin{enumerate}[resume]
  \item सवाल 6--10 एक शब्दकोश की ओर इशारा करते हैं: \omterm{def:g4:decimals:point}{दशमलव बिंदु} से
        ही दोहराने वाला $k$ अंकों का गुच्छा उसी गुच्छे को $k$ नौ के
        ऊपर लिखने के बराबर है ($0.727272\ldots = \frac{72}{99}$)।
        इस शब्दकोश से $0.037037\ldots$ की \omterm{def:g9:fractions:fraction}{भिन्न} बताओ, उसे सरल
        करो --- और चकित हो जाओ: कौन-सी हैरान कर देने वाली सरल
        \omterm{def:g9:fractions:fraction}{भिन्न} ``$037$'' का जाप करती है? भाग देकर जाँचो।
  \item शब्दकोश को एक सरकाव के साथ मिलाओ: $0.0272727\ldots$ को
        सबसे सरल \omterm{def:g9:fractions:fraction}{भिन्न} के रूप में लिखो।
  \item संख्या $0.101001000100001\ldots$ (एक $1$, फिर एक $0$, फिर
        एक $1$, फिर दो $0$, फिर एक $1$, फिर तीन $0$, और ऐसे ही
        आगे) कभी रुकती नहीं। क्या यह आवर्ती दशमलव है? तब भाग I उसके
        बारे में क्या कहता है --- क्या वह कोई \omterm{def:g9:fractions:fraction}{भिन्न} हो सकती है?
        (तुम्हारी पहली सिद्ध \emph{अपरिमेय} संख्या से मुलाक़ात हो
        चुकी; सबसे मशहूर अपरिमेय \cref{ch:g9:sqrt} में पकड़ी जाती
        है।)
  \item $2^{30}$ में मोटे तौर पर कितने अंक हैं?
        $2^{10} = 1\,024 \approx 1.024 \times 10^3$ और \omterm{def:g8:powers:def}{घातांक} के
        नियमों (\cref{thm:g9:fractions:rules}) से $2^{30}$ को
        (लगभग) \omterm{def:g9:fractions:scientific}{वैज्ञानिक लेखन} में लिखो, और नतीजा निकालो।
        (\cref{exo:g9:fractions:10} से तुलना करो।)
  \item आख़िरी दाँव: $0.20262026\ldots$ को \omterm{def:g9:fractions:fraction}{भिन्न} के रूप में लिखो,
        और पूरी समस्या की दोतरफ़ा सीख लिखो: कौन-से दशमलव लेखन
        भिन्नें हैं, और किन भिन्नों के कौन-से दशमलव लेखन होते हैं?
\end{enumerate}
\end{problem}
